% !TeX root = ../main.tex
\section{Цель работы}
Целью работы является экспериментальное исследование и анализ прогрешности позиционирования 
НАГ ГНСС в городских условиях

\section{Теоретическая часть}

\subsection{Сегменты СРНС}
Сетевая спутниковая радионавигационная система (СРНС) состоит из трёх основных сегментов:
\begin{itemize}
    \item \textbf{Космический сегмент:} орбитальная группировка навигационных спутников (НС) (штатно 24 КА в ГЛОНАСС). Спутники формируют радионавигационные сигналы и навигационные сообщения с использованием бортового синхронизирующего устройства (БСУ), задающего бортовую шкалу времени.
    \item \textbf{Сегмент контроля и управления (НКУ):} координационно-вычислительный центр (КВЦ), станции траекторных измерений (СТИ) и системный эталон времени и частоты. Задачи: эфемеридное обеспечение (прогноз орбит) и частотно-временное обеспечение (оценка уходов бортовых часов и передача частотно-временных поправок).
    \item \textbf{Сегмент потребителей (НАП):} приёмная аппаратура, выполняющая первичную обработку радиосигналов (оценка псевдозадержек $\tau$, доплеровских частот, декодирование данных) и вторичную обработку (расчёт пространственных координат и скорости).
\end{itemize}

\subsection{Псевдодальномерный метод позиционирования}
Шкала времени приёмника $t_r$ смещена относительно системной шкалы времени $t^s$ на величину $\Delta\tau(t_r)$. Измеряемая аппаратурой задержка сигнала является псевдозадержкой:
\begin{equation}
\tau(t_r) = t_r - t^s(t_r) = \tau_0(t_r) + \Delta\tau(t_r),
\end{equation}
где $\tau_0(t_r)$ --- истинное время распространения сигнала в пространстве. 

Псевдодальность до $i$-го спутника с координатами $(x_i, y_i, z_i)$ определяется как:
\begin{equation}
P_i = c \cdot \tau_i(t_r) = R_i + c\Delta\tau = \sqrt{(x - x_i)^2 + (y - y_i)^2 + (z - z_i)^2} + c\Delta\tau,
\end{equation}
где $c$ --- скорость света; $(x, y, z)$ --- искомые координаты потребителя.

Так как вектор неизвестных содержит 4 параметра $\mathbf{X} = [x, y, z, c\Delta\tau]^T$, для решения навигационной задачи требуется одновременный приём сигналов минимум четырёх видимых спутников ($N \ge 4$):
\begin{equation}
P_i = \sqrt{(x - x_i)^2 + (y - y_i)^2 + (z - z_i)^2} + c\Delta\tau, \quad i = 1, \dots, N.
\end{equation}
Система уравнений линеаризуется и решается относительно поправок к координатам методом наименьших квадратов.

\subsection{Системы координат и методы решения}
\begin{itemize}
    \item \textbf{Системы координат:} в СРНС расчёты ведутся в геоцентрических гринвичских прямоугольных системах ECEF (ПЗ-90.11 в ГЛОНАСС, WGS-84 в GPS), которые затем переводятся в геодезические координаты (широта $B$, долгота $L$, высота $H$ относительно референц-эллипсоида) или локальную систему координат ENU (East-North-Up).
    \item \textbf{Режимы позиционирования:} абсолютное кодовое позиционирование (SPP, точность $2\dots5\text{ м}$); дифференциальные измерения по коду (DGNSS); высокоточные фазовые режимы с сантиметровым разрешением (RTK/PPK) и режим высокоточного абсолютного позиционирования (PPP).
\end{itemize}
